% methods.tex -- input by main.tex
\section*{Materials and methods}

\subsection*{R1: Dual-FBS Isolation practice corpus and institutional recipes}
We assembled peer-reviewed Isolation or primary diagnostic isolation methods that state both a numeric pre-inoculation (growth) serum percentage and a numeric post-inoculation (maintenance or virus-culture) serum percentage, including explicit 0\%. Rows were coded only when both values were explicit in the source. Every dual numeric pair was retained with equal weight, including the common 10\%$\rightarrow$2\% pattern and less common equal holds, increases, or near-serum-free maintenance. The working table is \texttt{r1\_isolation\_standards\_and\_practice/expansion/isolation-refs-dual-fbs\_only.csv} ($n=170$ rows; search logs in \texttt{expansion/}).

Separately, we recorded published institutional recipes prescribing growth versus maintenance serum for diagnostic Isolation (ATCC; CLSI M41-A; ASM; NEADL/CIRAD PPRV SOP; WOAH Terrestrial Manual; WHO EPI; CDC measles laboratory tools; Health Canada/CCDR B95-a support; \texttt{institution-guideline-protocol-refs.md}). Eight guidelines with parseable dual percentages form the plotted overlay. USDA APHIS VIRPRO1013 is retained in the table but excluded from plots because it addresses master-seed extraneous-agent testing rather than diagnostic Isolation. Distributions were plotted with \texttt{expansion/plot\_fbs\_distributions.py}.

\emph{Public-records request.} To test whether reference and diagnostic laboratories follow the published step-down in practice, a public-records request was filed for the culture protocols in use (serum percentage in growth and maintenance media, and controls) (\texttt{r1\_isolation\_standards\_and\_practice/wrk/foia\_status.md}). \todoA{Agency, date filed, scope of request, number of protocols requested.} Released protocols will be coded with the same dual-FBS rule as the literature corpus.

\subsection*{R2: Isolation-stage negative controls}
Each of the 170 papers in the R1 dual-FBS corpus, plus three Isolation papers scored from the earlier unfiltered overview (all Tier D; 173 rows in total), was scored for the presence and quality of an Isolation-stage negative control on a four-tier scheme (\texttt{r2\_negative\_controls/control-cultures.md}). Tier A: an uninoculated culture handled as the Isolation condition, including the matched maintenance serum. Tier B: an uninoculated Isolation-stage control whose medium match is not stated. Tier C: a control present only at another stage (titration, plaque assay, immunofluorescence). Tier D: no qualifying control reported or full text unavailable. Ambiguity was resolved to the lower tier. Scores and quote-backed justifications are in \texttt{control-cultures-overview.csv} and \texttt{control-cultures-justification-notes.md}; three scoring passes are logged in \texttt{r2\_rescore\_status\_2026-09-19*.md}.

\subsection*{R3: Serum-only culture experiment and morphology labelling}
\emph{Culture experiment.} A U.S.-based contract research organisation (CRO) cultured Vero E6 cells (ATCC CRL-1586) with no inoculum along two identical paths that differed only in FBS concentration during the experimental stage: Path 1 (Culture A) at 10\% and Path 2 (Culture B) at 2\%. After thawing and three preparatory passages, passage-4 cultures were imaged daily for five days by phase-contrast light microscopy (EVOS FL, 10$\times$/20$\times$); transmission electron micrographs (Tecnai G2 Spirit BioTWIN) were also collected. Raw images are deposited on Zenodo (doi:10.5281/zenodo.17928456). \todoA{Medium base, seeding density, feeding schedule and TEM sampling days from the CRO report.}

\emph{Descriptor set.} CPE descriptors were extracted from Table 7 of CLSI M41-A (Table~\ref{tab:clsi-m41-cpe}): rounded, ballooned/enlarged, syncytia, vacuolation, detachment, granularity, refractile, cytoplasmic strands, nonspecific degeneration, no CPE. Mapping the CRO image descriptions onto this set yielded the six-term CRO-comparable set used for per-frame comparison: dying cells (Dy), rounded (Ro), vacuolation (V), detached (D), granularity (G) and refractile (Re). Healthy-type terms (e.g. look healthy, well-defined nuclei, cytoplasmic extensions) were coded from the same descriptions.

\begin{table}[htbp]
\centering
\small
\caption{Incidence table for CPE descriptors in Table 7 of CLSI M41}
\label{tab:clsi-m41-cpe}

\begin{tabular}{l *{10}{>{\centering\arraybackslash}m{0.5cm}} p{3.2cm}}
\toprule
\textbf{Virus}\rule{0pt}{4.8cm}
& \rotatebox[origin=br]{90}{\makebox[0pt][l]{\textbf{Rounded}}}
& \rotatebox[origin=br]{90}{\makebox[0pt][l]{\textbf{Ballooned/Enlarged}}}
& \rotatebox[origin=br]{90}{\makebox[0pt][l]{\textbf{Syncytia}}}
& \rotatebox[origin=br]{90}{\makebox[0pt][l]{\textbf{Vacuolation}}}
& \rotatebox[origin=br]{90}{\makebox[0pt][l]{\textbf{Detached}}}
& \rotatebox[origin=br]{90}{\makebox[0pt][l]{\textbf{Granularity}}}
& \rotatebox[origin=br]{90}{\makebox[0pt][l]{\textbf{Refractile}}}
& \rotatebox[origin=br]{90}{\makebox[0pt][l]{\textbf{Cytoplasmic strands}}}
& \rotatebox[origin=br]{90}{\makebox[0pt][l]{\textbf{Nonspecific degeneration}}}
& \rotatebox[origin=br]{90}{\makebox[0pt][l]{\textbf{No CPE**}}} \\
\midrule
\arrayrulecolor{gray!30}
Adenovirus          & x & x &   &   &   &   &   &   &   &   \\
\hline
Cytomegalovirus     & x & x &   &   &   &   &   &   &   &   \\
\hline
Enterovirus         & x &   &   &   &   &   & x &   &   &   \\
\hline
Herpes Simplex      & x & x & x &   &   &   &   &   &   &   \\
\hline
Influenza           &   &   &   & x &   & x &   &   & x & x \\
\hline
Measles             &   & x & x & x &   & x &   &   &   &   \\
\hline
Metapneumovirus     & x &   & x &   & x &   & x &   &   &   \\
\hline
Mumps               & x &   & x &   &   & x &   &   & x &   \\
\hline
Parainfluenza viruses & x &   & x &   &   & x &   &   & x &   \\
\hline
Reoviruses          &   &   &   &   & x &   &   &   & x &   \\
\hline
Respiratory Syncytial Virus &   &   & x &   &   & x &   &   & x &   \\
\hline
Rhinoviruses        & x &   &   &   &   &   & x &   &   &   \\
\hline
SARS-CoV            & x &   &   &   & x &   & x &   &   &   \\
\hline
Vaccinia            & x & x & x &   &   &   & x & x &   &   \\
\hline
Varicella-zoster virus & x & x & x &   &   & x & x & x &   &   \\
\arrayrulecolor{black}
\bottomrule
\multicolumn{11}{l}{\textbf{**} (i.e. CLSI mentions no CPE in inoculated culture)} \\
\end{tabular}

\end{table}

\emph{CRO labels.} The CRO supplied written descriptions for groups of ten frames and for 22 individual frames (9 Culture A, 13 Culture B). Labels were coded from the report text by a fixed rule set (\texttt{r3\_image\_labeling/cro-results/build\_cro\_labels.py}); only a frame's own single-frame description contributes to per-frame analysis (coding record: \texttt{cro\_labels\_consolidation.md}).

\emph{Automated labelling.} Computer-vision tools benchmarked against the 22 CRO-described frames did not reach accuracy adequate to label the remaining frames (Zenodo record 19607215), so human raters were used.

\emph{Blinded independent raters.} One hundred passage-4 frames (50 per culture; days 1--5, ten fields per day) were renamed with blinded identifiers (\texttt{images\_rater\_blinded\_mapping.csv}) and rated by two independent researchers with routine phase-contrast cell-culture experience (IR1, IR2). Raters were told the images would be used to train and evaluate AI vision systems and were not told the culture conditions, the absence of inoculum, or the hypothesis. Pass 1 was free-text description; Pass 2 was a ten-item structured checklist (three healthy items, seven CPE items) graded Yes/Partial/Mild/No. Pass-1 text was coded with fixed regular-expression lexicons; Pass-2 marks were scored as strict incidence and as a graded score (Yes 1.0, Mild 0.66, Partial 0.5, Minimal 0.25, No/minimal 0.15, No 0). Confluence was not used as an indicator of culture health. Scripts and outputs: \texttt{r3\_image\_labeling/ir-results/} and \texttt{wrk/viz/}. One Culture B day-4 field appears twice in the set: the CRO originals for path~2 frames 407 and 408 are the same file (identical checksum), so every rater described that field twice. The duplicate (frame 408, blinded \texttt{CultureB\_day4\_08}) is excluded from all analyses for every rater (IR1, IR2, IR3v and CRO) through a documented exclusion list applied by the processing scripts (\texttt{r3\_image\_labeling/frames\_excluded\_from\_analysis.csv}); raw rater files are unchanged. All per-frame results therefore use 99 unique frames (50 Culture A, 49 Culture B). Frame 408 has no individual CRO description, so the 22 CRO-described frames are unaffected.

\emph{Briefing-condition pilot arm (IR3v).} A third rater (IR3v) rated the same 100 frames (99 unique frames analysed) under a briefing that differed from the IR1/IR2 brief. \todoA{Describe the briefing condition only after the IR3v rater has completed Pass 2 and been debriefed (anonymity and debrief guardrail).} IR3v Pass 1 text was scored with the unchanged shared lexicon and, separately, with a clarified coding that applies the rater's own answers to five clarification questions (\texttt{r3\_image\_labeling/ir-results/IR3v\_pass1\_side\_by\_side.md}): shrinking was coded as the CPE class rounded/shrinking (Ro); the misspelling ``vaculation'' as vacuolation (V), while sentences describing a single large bright vacuole in about 1\% of cells, which the rater distinguished from vacuolation, were coded separately; refraction/refractive as refractile (Re). For projections the rater gave a two-part rule (thin and late in the time course = stressed; thick and early = healthy-type); no description stated projection thickness, so projections were left uncoded. \pending{IR3v Pass 2 checklist in progress.}

\emph{Statistics.} Per-item agreement between raters is reported as percent agreement and Cohen's $\kappa$ where both raters vary. \todoA{Choose an inferential test for the A--B contrast (e.g. frame-level permutation test or mixed model with day as a factor); currently descriptive only.}

\subsection*{Ethics and participant anonymity}
Raters and the contract laboratory are identified only by role codes (CRO, IR1, IR2, IR3v). Raters were paid for their work and were not told the study hypothesis. \todoA{Ethics review / exemption statement as required by PLOS.}
